\section{线性结构专项章正文案例推导补强}

这一节补齐本章正文出现过的 LeetCode 案例推导。阅读顺序统一按“题目说明、样例入口、暴力基线、暴力过程、重复工作、优化条件、相邻状态、状态复用、指针和字段、代码落点、正确性、边界实验、复杂度变化、迁移追问”展开；目的不是新增题量，而是把正文中的案例都讲成可复述、可证明、可迁移的解题链。

\subsection{LeetCode 19 删除链表的倒数第 N 个结点}

\textbf{题目说明。} LeetCode 19 删除链表的倒数第 N 个结点 的题面入口要先还原成具体任务：倒数位置可以转成两个指针之间固定 n 步的距离。

\textbf{样例入口。} 拿长度 5 删除倒数第 2 个手算。快指针先走 2 步，慢指针和快指针保持固定距离；当快指针到尾后，慢指针正好停在待删节点前驱。再试删除头节点，若没有虚拟头就缺前驱。

\textbf{暴力基线。} 暴力先扫描一遍链表求长度 len，再计算要删除的是第 len-n+1 个节点；第二遍走到它的前驱后改 next。

\textbf{暴力过程。} 以 1->2->3->4->5、n=2 为例，暴力先数出长度 5，目标是正数第 4 个节点 4，再从头走到节点 3，把 3->next 改成 5。

\textbf{重复工作。} 题面模型：倒数位置可以转成两个指针之间固定 n 步的距离。暴力版的慢点要落到本题：浪费在为了定位倒数位置先完整数长度。倒数第 n 个节点和它前驱，可以通过快慢指针之间固定 n+1 步距离一次扫描定位。后面的优化只消掉这类重复工作，不能改变答案集合。

\textbf{优化条件。} 本题的关键观察是：快指针先走 n 步，再让快慢指针同步移动，慢指针停在待删节点前驱。这句话必须能翻译成可维护状态；如果题目条件改变后观察不再成立，就不能继续套原来的写法。

\textbf{相邻状态。} 规律是虚拟头统一头部删除，快慢指针固定距离定位前驱。把这个特例继续拆开看：每个指针都要有角色，上一段尾、当前段头、当前节点和后继入口不能混。只要一次改 next 之前没有保存后继，就可能丢掉整段未处理链。

\textbf{状态复用。} 使用虚拟头 dummy 指向 head。fast 先从 dummy 走 n+1 步，然后 fast 和 slow 同步走；fast 到空时，slow 正好是待删节点前驱。手算时只改被新输入影响的那一小部分，而不是回头重算完整候选。

\textbf{指针和字段。} 状态字段要能说清：dummy 负责统一删除头节点，fast 负责制造距离，slow 停在前驱，slow->next 是要删除的节点。手算时按这个协议跟踪：删除 1->2->3->4->5 的倒数第 2 个时，fast 先领先 slow 三步；同步移动到 fast 为空，slow 在 3，删除 slow->next=4。

\textbf{代码落点。} 关键代码是 dummy.next=head，fast=\&dummy，预走 n+1 步；最后 slow->next=slow->next->next。预走从 dummy 开始，才能覆盖删除头节点。关键代码行必须能逐句翻译成“旧状态怎样变成新状态”，否则就只是模板。

\textbf{正确性。} fast 和 slow 始终相隔 n+1 条边；当 fast 走到链表末尾之后，slow 后面正好还有 n 个节点，所以 slow 是目标前驱。

\textbf{边界实验。} 先用删除头节点、n 等于链表长度、只有一个节点、题目保证 n 合法做反例检查。实验时覆盖删除头节点、n 等于链表长度、只有一个节点、题目保证 n 合法，每步画出前驱、当前、后继和下一段头。链表题不要只看输出值，还要检查节点是否丢失或成环。

\textbf{复杂度变化。} 两遍扫描 O(n)；快慢指针仍是 O(n) 时间，但只走一遍主体扫描，额外空间 O(1)。

\textbf{迁移追问。} 如果题目改成“若 n 不保证合法，需要在快指针预走阶段检测长度不足”，先判断原状态是否仍然覆盖新答案集合，再决定是微调边界、改 value 语义，还是换算法。

\subsection{LeetCode 21 合并两个有序链表}

\textbf{题目说明。} LeetCode 21 合并两个有序链表 的题面入口要先还原成具体任务：两个链表已经各自有序，下一节点只能来自两个当前头中较小者。

\textbf{样例入口。} 拿 1->3 和 2->4 手算，先接 1，再比较 3 和 2 接 2，之后接 3 和 4。

\textbf{暴力基线。} 暴力可以把两个链表的值全部放入数组，排序后重新建链表。

\textbf{暴力过程。} 以 1->3 和 2->4 为例，数组法得到 [1, 2, 3, 4]；但原链表本来有序，每一步只需要比较两个当前头。

\textbf{重复工作。} 题面模型：两个链表已经各自有序，下一节点只能来自两个当前头中较小者。暴力版的慢点要落到本题：浪费在完整排序和新建节点。两个链表各自有序，下一节点一定来自 l1 或 l2 的较小头节点。后面的优化只消掉这类重复工作，不能改变答案集合。

\textbf{优化条件。} 本题的关键观察是：虚拟头统一结果链表，尾指针不断接上较小节点并推进来源链。这句话必须能翻译成可维护状态；如果题目条件改变后观察不再成立，就不能继续套原来的写法。

\textbf{相邻状态。} 规律是结果尾指针每次接较小头节点，某条链耗尽后直接接另一条剩余部分；虚拟头统一处理首节点。把这个特例继续拆开看：每个指针都要有角色，上一段尾、当前段头、当前节点和后继入口不能混。只要一次改 next 之前没有保存后继，就可能丢掉整段未处理链。

\textbf{状态复用。} 用 dummy 作为结果虚拟头，tail 指向结果尾。每次比较 l1 和 l2，接上较小节点并推进对应链表。手算时只改被新输入影响的那一小部分，而不是回头重算完整候选。

\textbf{指针和字段。} 状态字段要能说清：l1/l2 是两条未处理链的当前头，tail 是已合并链表尾部，dummy.next 是最终结果入口。手算时按这个协议跟踪：1 和 2 比较接 1；3 和 2 比较接 2；3 和 4 比较接 3；最后 l1 空，直接接剩余 4。

\textbf{代码落点。} 关键是 tail->next=chosen 后 tail=tail->next，再推进 chosen 来源链。某条链为空后 tail->next 接另一条剩余链。关键代码行必须能逐句翻译成“旧状态怎样变成新状态”，否则就只是模板。

\textbf{正确性。} 两条链都已排序，两个当前头中较小者不可能被另一条链后面的节点超过，所以它就是全局下一个节点。

\textbf{边界实验。} 先用某一条链为空、重复值、剩余整段接回、是否复用原节点做反例检查。实验时覆盖某一条链为空、重复值、剩余整段接回、是否复用原节点，每步画出前驱、当前、后继和下一段头。链表题不要只看输出值，还要检查节点是否丢失或成环。

\textbf{复杂度变化。} 排序数组法 O((m+n)log(m+n)) 且额外空间高；双指针合并 O(m+n) 时间，O(1) 额外空间。

\textbf{迁移追问。} 如果题目改成“合并 K 条链表可用分治或堆，核心仍是维护每条链的当前头”，先判断原状态是否仍然覆盖新答案集合，再决定是微调边界、改 value 语义，还是换算法。

\subsection{LeetCode 25 K 个一组翻转链表}

\textbf{题目说明。} LeetCode 25 K 个一组翻转链表 的题面入口要先还原成具体任务：链表被切成长度为 K 的连续组，只有完整组才反转。

\textbf{样例入口。} 拿 1->2->3->4、k=2 画图。第一组边界是 1 和 2，下一组入口是 3；反转前如果不保存 3，链会断。反转后上一组尾接 2，旧头 1 接 3。

\textbf{暴力基线。} 暴力可以把每 K 个节点的值放进数组反转，再写回链表；这避开了指针重连，但没有真正按节点反转。

\textbf{暴力过程。} 以 1->2->3->4->5、k=2 为例，第一组 1、2 反转成 2->1，第二组 3、4 反转成 4->3，剩余 5 不动。

\textbf{重复工作。} 题面模型：链表被切成长度为 K 的连续组，只有完整组才反转。暴力版的慢点要落到本题：浪费在数组辅助和节点值交换。题目要求链表节点重排时，应直接反转每一组的 next 指针。后面的优化只消掉这类重复工作，不能改变答案集合。

\textbf{优化条件。} 本题的关键观察是：每轮先探测 K 个节点确认完整，再保存下一组头并做局部反转接回。这句话必须能翻译成可维护状态；如果题目条件改变后观察不再成立，就不能继续套原来的写法。

\textbf{相邻状态。} 规律是每组先找 kth 和 group\_next，再局部反转并接回前后两段。把这个特例继续拆开看：每个指针都要有角色，上一段尾、当前段头、当前节点和后继入口不能混。只要一次改 next 之前没有保存后继，就可能丢掉整段未处理链。

\textbf{状态复用。} dummy 指向头。group\_prev 是上一组尾，先从 group\_prev 出发找第 k 个节点 group\_end；不足 k 个就停止。保存 next\_group 后反转 [group\_start, group\_end]，再接回。手算时只改被新输入影响的那一小部分，而不是回头重算完整候选。

\textbf{指针和字段。} 状态字段要能说清：group\_prev 是已处理部分尾，group\_start 是本组头，group\_end 是本组尾，next\_group 是下一组入口。手算时按这个协议跟踪：第一组 start=1、end=2、next\_group=3。反转后得到 2->1，再让 group\_prev->next=2，1->next=3。

\textbf{代码落点。} 关键顺序是先保存 next\_group，再局部反转；反转结束后原 group\_start 变成本组尾，下一轮 group\_prev=group\_start。关键代码行必须能逐句翻译成“旧状态怎样变成新状态”，否则就只是模板。

\textbf{正确性。} 每次只改变完整 K 组内部指针，并把本组头尾接回已处理链和未处理链；不足 K 个节点不进入反转，符合题意。

\textbf{边界实验。} 先用不足 K 个不反转、K 为一、刚好整除、反转后头尾接错做反例检查。实验时覆盖不足 K 个不反转、K 为一、刚好整除、反转后头尾接错，每步画出前驱、当前、后继和下一段头。链表题不要只看输出值，还要检查节点是否丢失或成环。

\textbf{复杂度变化。} 数组辅助 O(n) 空间；原地分组反转每节点访问常数次，O(n) 时间、O(1) 空间。

\textbf{迁移追问。} 如果题目改成“递归写法短但可能有栈风险，迭代写法更接近工程实现”，先判断原状态是否仍然覆盖新答案集合，再决定是微调边界、改 value 语义，还是换算法。

\subsection{LeetCode 155 最小栈}

\textbf{题目说明。} LeetCode 155 最小栈 的题面入口要先还原成具体任务：普通栈只能看到栈顶，不能常数时间知道历史最小值。

\textbf{样例入口。} 拿 push 2、push 0、push 3、push 0 手算，当前最小值序列是 2、0、0、0；弹出最后一个 0 后，最小值仍是 0。

\textbf{暴力基线。} 慢版本让每个位置向左或向右线性寻找边界，或者让每个结构反复等待匹配。它能算出答案，但很多尚未完成的候选会被未来同一个事件一起解决。放到本题就是：先按“普通栈只能看到栈顶，不能常数时间知道历史最小值”完整试慢版本；样例里已经能看到“规律是辅助栈要在每个深度同步保存当前最小值，重复最小值不能只保存一次”，所以优化前必须先能说清慢版本枚举了哪些候选。

\textbf{暴力过程。} 拿 push 2、push 0、push 3、push 0 手算，当前最小值序列是 2、0、0、0；弹出最后一个 0 后，最小值仍是 0。

\textbf{重复工作。} 题面模型：普通栈只能看到栈顶，不能常数时间知道历史最小值。暴力版的慢点要落到本题：慢版本会围绕“普通栈只能看到栈顶，不能常数时间知道历史最小值”反复寻找最近边界或最近未完成对象。样例规律是“规律是辅助栈要在每个深度同步保存当前最小值，重复最小值不能只保存一次”，说明旧候选一旦遇到当前输入就能被批量结算；继续为每个位置单独向左或向右扫描就是浪费。后面的优化只消掉这类重复工作，不能改变答案集合。

\textbf{优化条件。} 本题的关键观察是：辅助栈同步保存每个深度对应的最小值。这句话必须能翻译成可维护状态；如果题目条件改变后观察不再成立，就不能继续套原来的写法。

\textbf{相邻状态。} 规律是辅助栈要在每个深度同步保存当前最小值，重复最小值不能只保存一次。把这个特例继续拆开看：栈里的对象都是答案尚未确定的候选；一旦当前元素触发弹栈，被弹出的候选必须已经同时拿到了左边界和右边界，未来元素不再有资格改变它的答案。

\textbf{状态复用。} 把反复寻找最近边界改成维护未完成候选；新元素只和栈顶比较，连续弹出一批已经被当前元素解决的对象。本题落点是：规律是辅助栈要在每个深度同步保存当前最小值，重复最小值不能只保存一次。手算时只改被新输入影响的那一小部分，而不是回头重算完整候选。

\textbf{指针和字段。} 状态字段要能说清：栈内对象的业务含义、当前输入、弹出时结算的对象、左边界和右边界；栈中存下标还是值要明确。手算时按这个协议跟踪：拿 push 2、push 0、push 3、push 0 手算，当前最小值序列是 2、0、0、0；弹出最后一个 0 后，最小值仍是 0。然后按协议复查：手算时记录栈内元素的业务含义，不只记录数值。入栈表示答案未定，弹栈表示当前事件已经让它的答案定下来。

\textbf{代码落点。} 弹栈前确认栈非空，弹出后再读取新的左边界；相等元素和哨兵高度要有统一策略。关键代码行必须能逐句翻译成“旧状态怎样变成新状态”，否则就只是模板。

\textbf{正确性。} 证明从弹出时机开始：被弹出的对象在当前时刻答案已经确定，未来元素无法再改变它。

\textbf{边界实验。} 先用重复最小值、弹出最小值、空栈操作约定做反例检查。实验时准备重复最小值、弹出最小值、空栈操作约定，打印每次入栈、弹栈和结算对象。重点观察相等元素、空栈、哨兵和最后一次结算。

\textbf{复杂度变化。} 暴力为每个位置寻找边界常见为平方级；单调栈让每个元素最多入栈一次、出栈一次，因此主体扫描为线性时间。

\textbf{迁移追问。} 如果题目改成“也可保存差值压缩空间，但可读性和溢出处理更复杂”，先判断原状态是否仍然覆盖新答案集合，再决定是微调边界、改 value 语义，还是换算法。

\subsection{LeetCode 84 柱状图中最大的矩形}

\textbf{题目说明。} LeetCode 84 柱状图中最大的矩形 的题面入口要先还原成具体任务：每根柱子作为最矮高度时，需要知道左右第一个更矮位置。

\textbf{样例入口。} 拿高度 [2, 1, 2] 手算。第一个 2 遇到更矮的 1 时，右边界已经确定在当前位置，左边界是弹出后新栈顶的后一位，所以面积可结算。若一直递增，最后需要哨兵触发结算。

\textbf{暴力基线。} 暴力枚举每根柱子作为矩形最矮高度，再向左右扩展直到遇到更矮柱子，计算最大面积。

\textbf{暴力过程。} 以 heights=[2, 1, 5, 6, 2, 3] 为例，高度 5 的柱子向右可覆盖 6，向左遇到 1 停止，面积可为 5*2=10。

\textbf{重复工作。} 题面模型：每根柱子作为最矮高度时，需要知道左右第一个更矮位置。暴力版的慢点要落到本题：浪费在每根柱子都重新找左右第一个更矮位置。单调递增栈能在遇到更矮柱时一次结算被弹出柱子的右边界。后面的优化只消掉这类重复工作，不能改变答案集合。

\textbf{优化条件。} 本题的关键观察是：单调递增栈在遇到更矮柱时结算被弹出柱子的最大宽度。这句话必须能翻译成可维护状态；如果题目条件改变后观察不再成立，就不能继续套原来的写法。

\textbf{相邻状态。} 规律是栈保存递增高度下标，弹出时才知道该柱子的最大扩展宽度。把这个特例继续拆开看：栈里的对象都是答案尚未确定的候选；一旦当前元素触发弹栈，被弹出的候选必须已经同时拿到了左边界和右边界，未来元素不再有资格改变它的答案。

\textbf{状态复用。} 栈保存下标，且高度单调递增。扫描到更矮高度时，弹出 top，当前下标是右边界，弹出后新的栈顶是左边界。手算时只改被新输入影响的那一小部分，而不是回头重算完整候选。

\textbf{指针和字段。} 状态字段要能说清：top 是被结算柱子，height[top] 是矩形高度，left\_smaller=stack.top，right\_smaller=i，宽度为 i-left\_smaller-1。手算时按这个协议跟踪：扫描到高度 2 时，6 和 5 依次被弹出；弹出 5 时左边界是高度 1 的下标，右边界是当前 2 的下标，面积 10。

\textbf{代码落点。} 关键是两端加哨兵 0，或循环结束后手动清栈。相等高度要有统一策略，常见写法用 >= 弹出保证边界一致。关键代码行必须能逐句翻译成“旧状态怎样变成新状态”，否则就只是模板。

\textbf{正确性。} 柱子被弹出时，左右第一个更矮位置都已经确定，未来更右的柱子不能再扩大以它为最低高度的矩形。

\textbf{边界实验。} 先用递增、递减、相等高度、哨兵零、宽度计算做反例检查。实验时准备递增、递减、相等高度、哨兵零、宽度计算，打印每次入栈、弹栈和结算对象。重点观察相等元素、空栈、哨兵和最后一次结算。

\textbf{复杂度变化。} 暴力 O(\ensuremath{n^{2}})，单调栈每个下标入栈出栈一次，O(n) 时间、O(n) 空间。

\textbf{迁移追问。} 如果题目改成“最大矩形可以作为二维矩阵最大矩形的行压缩子问题”，先判断原状态是否仍然覆盖新答案集合，再决定是微调边界、改 value 语义，还是换算法。

\subsection{LeetCode 239 滑动窗口最大值}

\textbf{题目说明。} LeetCode 239 滑动窗口最大值 的题面入口要先还原成具体任务：每个窗口最大值来自仍未过期且没有被新元素支配的候选。

\textbf{样例入口。} 拿 [1, 3, -1]、窗口 3 手算，队尾的 1 在 3 进入后永远不可能成为最大值，因为 3 更新且更靠右。于是 1 可以弹出。

\textbf{暴力基线。} 暴力枚举每个长度为 k 的窗口，扫描窗口内所有元素求最大值。

\textbf{暴力过程。} 以 [1, 3, -1, -3, 5]、k=3 为例，第一个窗口 [1, 3, -1] 最大值为 3；下一个窗口 [3, -1, -3] 仍要重新扫描会重复看 3 和 -1。

\textbf{重复工作。} 题面模型：每个窗口最大值来自仍未过期且没有被新元素支配的候选。暴力版的慢点要落到本题：浪费在相邻窗口重复求最大值。若新来的数更大且位置更靠右，队尾所有不大于它的旧数未来都不可能再当最大值。后面的优化只消掉这类重复工作，不能改变答案集合。

\textbf{优化条件。} 本题的关键观察是：单调队列保存下标，队尾小于等于新值的候选被弹出。这句话必须能翻译成可维护状态；如果题目条件改变后观察不再成立，就不能继续套原来的写法。

\textbf{相邻状态。} 规律是单调队列保存可能成为当前或未来窗口最大值的下标，过期下标从队首清理。把这个特例继续拆开看：栈里的对象都是答案尚未确定的候选；一旦当前元素触发弹栈，被弹出的候选必须已经同时拿到了左边界和右边界，未来元素不再有资格改变它的答案。

\textbf{状态复用。} 单调递减队列保存下标。入队前从队尾弹出 nums[back]<=nums[i] 的下标；队首若过期 i-k 则弹出；窗口形成后队首就是最大值。手算时只改被新输入影响的那一小部分，而不是回头重算完整候选。

\textbf{指针和字段。} 状态字段要能说清：deque.front 是当前最大值下标，i 是右端，i-k+1 是当前窗口左端。手算时按这个协议跟踪：1 入队后，3 进入时弹出 1，因为 3 更大且更晚；窗口形成后队首 3 给出最大值。

\textbf{代码落点。} 关键是队列保存下标而不是值，这样才能判断过期。顺序通常是清队尾、入队、清队首过期、读答案。关键代码行必须能逐句翻译成“旧状态怎样变成新状态”，否则就只是模板。

\textbf{正确性。} 被队尾新元素弹出的旧元素既更小又更早，在未来任何包含旧元素的窗口里也会包含新元素，因此不会成为最大值。

\textbf{边界实验。} 先用k 为一、k 等于数组长度、重复最大值、队首过期做反例检查。实验时准备k 为一、k 等于数组长度、重复最大值、队首过期，打印每次入栈、弹栈和结算对象。重点观察相等元素、空栈、哨兵和最后一次结算。

\textbf{复杂度变化。} 暴力 O(nk)，单调队列每个下标入队出队一次，O(n) 时间、O(k) 空间。

\textbf{迁移追问。} 如果题目改成“受限子序列和使用同样队列维护最近 k 个 DP 最大值”，先判断原状态是否仍然覆盖新答案集合，再决定是微调边界、改 value 语义，还是换算法。

\subsection{LeetCode 141 环形链表}

\textbf{题目说明。} LeetCode 141 环形链表 的题面入口要先还原成具体任务：快慢指针在有环时必然相遇，无环时快指针先到空。

\textbf{样例入口。} 拿 1->2->3 且 3 指回 2 手算，快指针每次走两步，慢指针每次走一步，进入环后二者距离会被不断缩小直到相遇。

\textbf{暴力基线。} 慢版本可以先把节点顺序抄到数组里，按数组推导目标顺序。回到链表后，难点不再是值的顺序，而是节点身份和 next 指针不能丢。放到本题就是：先按“快慢指针在有环时必然相遇，无环时快指针先到空”完整试慢版本；样例里已经能看到“规律是快慢指针的相遇说明有环；无环时快指针会先到空”，所以优化前必须先能说清慢版本枚举了哪些候选。

\textbf{暴力过程。} 拿 1->2->3 且 3 指回 2 手算，快指针每次走两步，慢指针每次走一步，进入环后二者距离会被不断缩小直到相遇。

\textbf{重复工作。} 题面模型：快慢指针在有环时必然相遇，无环时快指针先到空。暴力版的慢点要落到本题：慢版本会围绕“快慢指针在有环时必然相遇，无环时快指针先到空”借助数组或多次扫描确认目标顺序。样例规律是“规律是快慢指针的相遇说明有环；无环时快指针会先到空”，说明优化点不在节点值，而在少量指针角色能否保持已处理链、未处理链和接回位置都不丢。后面的优化只消掉这类重复工作，不能改变答案集合。

\textbf{优化条件。} 本题的关键观察是：不用哈希表也能常数空间检测环。这句话必须能翻译成可维护状态；如果题目条件改变后观察不再成立，就不能继续套原来的写法。

\textbf{相邻状态。} 规律是快慢指针的相遇说明有环；无环时快指针会先到空。把这个特例继续拆开看：每个指针都要有角色，上一段尾、当前段头、当前节点和后继入口不能混。只要一次改 next 之前没有保存后继，就可能丢掉整段未处理链。

\textbf{状态复用。} 把数组辅助结构换成几个稳定指针角色；重连顺序必须保证已处理链合法、未处理链仍可达。本题落点是：规律是快慢指针的相遇说明有环；无环时快指针会先到空。手算时只改被新输入影响的那一小部分，而不是回头重算完整候选。

\textbf{指针和字段。} 状态字段要能说清：虚拟头、上一段尾、当前段头、当前节点、后继节点；任何改 next 的操作之前都要保存后续入口。手算时按这个协议跟踪：拿 1->2->3 且 3 指回 2 手算，快指针每次走两步，慢指针每次走一步，进入环后二者距离会被不断缩小直到相遇。然后按协议复查：手算时画出 prev、cur、next 和下一段头节点。每次改 next 前先保存后继，这是链表推导的安全线。

\textbf{代码落点。} 所有重连步骤按保存后继、断开或反转、接回前缀、接回后缀的顺序写；最后检查是否形成环或丢节点。关键代码行必须能逐句翻译成“旧状态怎样变成新状态”，否则就只是模板。

\textbf{正确性。} 证明从链结构开始：已处理链连通、未处理链可达、二者连接点明确，节点没有丢失也没有成环。

\textbf{边界实验。} 先用空链、单节点自环、两节点环、尾部无环做反例检查。实验时覆盖空链、单节点自环、两节点环、尾部无环，每步画出前驱、当前、后继和下一段头。链表题不要只看输出值，还要检查节点是否丢失或成环。

\textbf{复杂度变化。} 暴力可能多次扫描链表或拷贝到数组；优化后每个节点通常只被访问、重连常数次，目标是线性时间和常数额外空间。

\textbf{迁移追问。} 如果题目改成“若要求环入口，用相遇后双指针同步走”，先判断原状态是否仍然覆盖新答案集合，再决定是微调边界、改 value 语义，还是换算法。

\subsection{LeetCode 146 LRU 缓存}

\textbf{题目说明。} LeetCode 146 LRU 缓存 的题面入口要先还原成具体任务：哈希表负责按 key 定位，双向链表负责维护新旧顺序。

\textbf{样例入口。} 拿容量 2 依次 put(1)、put(2)、get(1)、put(3) 手算。get(1) 后 1 变成最新，淘汰时应删 2。数组维护顺序每次移动成本高；链表移动节点到头 O(1)，哈希表按 key 定位节点 O(1)。

\textbf{暴力基线。} 暴力用数组保存缓存项，每次 get 线性查找 key，命中后把元素移动到最新位置；put 时线性找 key 或淘汰最旧项。

\textbf{暴力过程。} 容量 2，put(1)、put(2)、get(1)、put(3) 时，get(1) 会让 1 变成最新，put(3) 应淘汰最旧的 2。

\textbf{重复工作。} 题面模型：哈希表负责按 key 定位，双向链表负责维护新旧顺序。暴力版的慢点要落到本题：浪费在按 key 查找和移动新旧顺序都线性。LRU 需要两个能力：O(1) 按 key 定位节点，O(1) 把节点移动到头部和删除尾部。后面的优化只消掉这类重复工作，不能改变答案集合。

\textbf{优化条件。} 本题的关键观察是：访问和更新都把节点移动到头部，容量满时删除尾部最旧节点。这句话必须能翻译成可维护状态；如果题目条件改变后观察不再成立，就不能继续套原来的写法。

\textbf{相邻状态。} 规律是哈希负责找节点，双向链表负责维护新旧顺序。把这个特例继续拆开看：每个指针都要有角色，上一段尾、当前段头、当前节点和后继入口不能混。只要一次改 next 之前没有保存后继，就可能丢掉整段未处理链。

\textbf{状态复用。} 哈希表 key->链表节点迭代器，双向链表按新到旧保存 key/value。访问或更新节点时移动到头部，容量超限删除尾部。手算时只改被新输入影响的那一小部分，而不是回头重算完整候选。

\textbf{指针和字段。} 状态字段要能说清：map 定位缓存项，list.front 是最新项，list.back 是最旧项，capacity 控制淘汰边界。手算时按这个协议跟踪：get(1) 命中后，把 1 对应节点移到表头；put(3) 超容量时，删除表尾 2，并从哈希表同步 erase。

\textbf{代码落点。} C++ 可用 std::list<pair<int, int>> 保存顺序，unordered\_map<int, list::iterator> 定位。更新已有 key 时先改值再 splice 到头部。关键代码行必须能逐句翻译成“旧状态怎样变成新状态”，否则就只是模板。

\textbf{正确性。} 每次访问都把对应项放到最新位置，所以链表尾部始终是最长时间未访问项；哈希表和链表同步维护同一批 key。

\textbf{边界实验。} 先用容量为一、重复 put、get 不存在、更新已有值做反例检查。实验时覆盖容量为一、重复 put、get 不存在、更新已有值，每步画出前驱、当前、后继和下一段头。链表题不要只看输出值，还要检查节点是否丢失或成环。

\textbf{复杂度变化。} 数组暴力 get/put O(n)；哈希表加双链表 get/put 平均 O(1)，空间 O(capacity)。

\textbf{迁移追问。} 如果题目改成“C++ 可用 \code{std::list} 加迭代器，也可手写双向链表理解所有权”，先判断原状态是否仍然覆盖新答案集合，再决定是微调边界、改 value 语义，还是换算法。

\subsection{LeetCode 42 接雨水}

\textbf{题目说明。} LeetCode 42 接雨水 的题面入口要先还原成具体任务：每个位置的水由左侧最高和右侧最高的较小者决定。

\textbf{样例入口。} 拿高度 [0, 2, 0, 3] 手算，中间的 0 左侧最高是 2，右侧最高是 3，所以能接 2 格水。若用单调栈，遇到 3 时弹出 0，左边界是 2，右边界是 3，宽度为 1，高度差为 2。

\textbf{暴力基线。} 暴力对每个位置 i，向左扫描最高柱 leftMax，向右扫描最高柱 rightMax，当前位置接水为 min(leftMax, rightMax)-height[i]。

\textbf{暴力过程。} 以 [0, 1, 0, 2] 为例，下标 2 左侧最高为 1，右侧最高为 2，所以这里能接 1 格水。

\textbf{重复工作。} 题面模型：每个位置的水由左侧最高和右侧最高的较小者决定。暴力版的慢点要落到本题：浪费在每个位置都重新向两边找最高柱。左右最高值可以前缀后缀预处理，也可以用双指针增量维护。后面的优化只消掉这类重复工作，不能改变答案集合。

\textbf{优化条件。} 本题的关键观察是：前后缀、双指针、单调栈都是不同状态维护方式；要能说明各自结算对象。这句话必须能翻译成可维护状态；如果题目条件改变后观察不再成立，就不能继续套原来的写法。

\textbf{相邻状态。} 规律是每个凹槽必须同时等到左右边界，弹栈那一刻才真正结算水量。把这个特例继续拆开看：栈里的对象都是答案尚未确定的候选；一旦当前元素触发弹栈，被弹出的候选必须已经同时拿到了左边界和右边界，未来元素不再有资格改变它的答案。

\textbf{状态复用。} 双指针维护 leftMax 和 rightMax。每次处理较低一侧，因为该侧的水量已经由它自己的最高值和另一侧不低的边界确定。手算时只改被新输入影响的那一小部分，而不是回头重算完整候选。

\textbf{指针和字段。} 状态字段要能说清：left/right 是未处理区两端，leftMax/rightMax 是两侧已见最高柱，water 累加已确定位置的贡献。手算时按这个协议跟踪：若 height[left] < height[right]，右侧至少存在 height[right] 这个边界，left 位置能接多少只取决于 leftMax。

\textbf{代码落点。} 关键是先更新 leftMax/rightMax，再累加 max(0, leftMax-height[left]) 或 max(0, rightMax-height[right])。移动的是较低一侧。关键代码行必须能逐句翻译成“旧状态怎样变成新状态”，否则就只是模板。

\textbf{正确性。} 较低一侧的对侧边界已经足够高，不会成为短板；未来内部柱子只能影响未处理位置，不会改变当前已结算位置的水量。

\textbf{边界实验。} 先用单调递增、单调递减、平台高度、多个凹槽、数组长度不足三做反例检查。实验时准备单调递增、单调递减、平台高度、多个凹槽、数组长度不足三，打印每次入栈、弹栈和结算对象。重点观察相等元素、空栈、哨兵和最后一次结算。

\textbf{复杂度变化。} 逐点双向扫描 O(\ensuremath{n^{2}})，前后缀 O(n) 空间，双指针 O(n) 时间、O(1) 空间。

\textbf{迁移追问。} 如果题目改成“若柱子变成二维网格，需要最小堆从边界向内扩展”，先判断原状态是否仍然覆盖新答案集合，再决定是微调边界、改 value 语义，还是换算法。

\subsection{LeetCode 227 基本计算器 II}

\textbf{题目说明。} LeetCode 227 基本计算器 II 的题面入口要先还原成具体任务：乘除优先级高于加减，可以在读到下一个运算符时结算上一项。

\textbf{样例入口。} 拿 3+2*2 手算，加号前的 3 可以先入栈，乘号遇到 2 时要立刻把栈顶 2 乘以下一个数 2，而不能等到最后按顺序相加。

\textbf{暴力基线。} 慢版本让每个位置向左或向右线性寻找边界，或者让每个结构反复等待匹配。它能算出答案，但很多尚未完成的候选会被未来同一个事件一起解决。放到本题就是：先按“乘除优先级高于加减，可以在读到下一个运算符时结算上一项”完整试慢版本；样例里已经能看到“规律是栈保存已确定符号的项，乘除立即修改栈顶，加减压入正负项”，所以优化前必须先能说清慢版本枚举了哪些候选。

\textbf{暴力过程。} 拿 3+2*2 手算，加号前的 3 可以先入栈，乘号遇到 2 时要立刻把栈顶 2 乘以下一个数 2，而不能等到最后按顺序相加。

\textbf{重复工作。} 题面模型：乘除优先级高于加减，可以在读到下一个运算符时结算上一项。暴力版的慢点要落到本题：慢版本会围绕“乘除优先级高于加减，可以在读到下一个运算符时结算上一项”反复寻找最近边界或最近未完成对象。样例规律是“规律是栈保存已确定符号的项，乘除立即修改栈顶，加减压入正负项”，说明旧候选一旦遇到当前输入就能被批量结算；继续为每个位置单独向左或向右扫描就是浪费。后面的优化只消掉这类重复工作，不能改变答案集合。

\textbf{优化条件。} 本题的关键观察是：栈保存已经确定符号的项，乘除直接修改栈顶，加减压入正负项。这句话必须能翻译成可维护状态；如果题目条件改变后观察不再成立，就不能继续套原来的写法。

\textbf{相邻状态。} 规律是栈保存已确定符号的项，乘除立即修改栈顶，加减压入正负项。把这个特例继续拆开看：栈里的对象都是答案尚未确定的候选；一旦当前元素触发弹栈，被弹出的候选必须已经同时拿到了左边界和右边界，未来元素不再有资格改变它的答案。

\textbf{状态复用。} 把反复寻找最近边界改成维护未完成候选；新元素只和栈顶比较，连续弹出一批已经被当前元素解决的对象。本题落点是：规律是栈保存已确定符号的项，乘除立即修改栈顶，加减压入正负项。手算时只改被新输入影响的那一小部分，而不是回头重算完整候选。

\textbf{指针和字段。} 状态字段要能说清：栈内对象的业务含义、当前输入、弹出时结算的对象、左边界和右边界；栈中存下标还是值要明确。手算时按这个协议跟踪：拿 3+2*2 手算，加号前的 3 可以先入栈，乘号遇到 2 时要立刻把栈顶 2 乘以下一个数 2，而不能等到最后按顺序相加。然后按协议复查：手算时记录栈内元素的业务含义，不只记录数值。入栈表示答案未定，弹栈表示当前事件已经让它的答案定下来。

\textbf{代码落点。} 弹栈前确认栈非空，弹出后再读取新的左边界；相等元素和哨兵高度要有统一策略。关键代码行必须能逐句翻译成“旧状态怎样变成新状态”，否则就只是模板。

\textbf{正确性。} 证明从弹出时机开始：被弹出的对象在当前时刻答案已经确定，未来元素无法再改变它。

\textbf{边界实验。} 先用多位数、空格、除法向零截断、最后一个数字需要结算做反例检查。实验时准备多位数、空格、除法向零截断、最后一个数字需要结算，打印每次入栈、弹栈和结算对象。重点观察相等元素、空栈、哨兵和最后一次结算。

\textbf{复杂度变化。} 暴力为每个位置寻找边界常见为平方级；单调栈让每个元素最多入栈一次、出栈一次，因此主体扫描为线性时间。

\textbf{迁移追问。} 如果题目改成“若加入括号，需要再引入上下文栈或递归下降解析”，先判断原状态是否仍然覆盖新答案集合，再决定是微调边界、改 value 语义，还是换算法。

\subsection{LeetCode 160 相交链表}

\textbf{题目说明。} LeetCode 160 相交链表 的题面入口要先还原成具体任务：相交是节点地址相同，不是节点值相同。

\textbf{样例入口。} 拿 A: a1->a2->c1->c2 和 B: b1->c1->c2 手算，交点是节点地址 c1，不是值相等。两个指针走到尾后切换到对方头部，会在总路程相同后于 c1 相遇。

\textbf{暴力基线。} 慢版本可以先把节点顺序抄到数组里，按数组推导目标顺序。回到链表后，难点不再是值的顺序，而是节点身份和 next 指针不能丢。放到本题就是：先按“相交是节点地址相同，不是节点值相同”完整试慢版本；样例里已经能看到“规律是比较节点身份，长度差由切换链表自动抵消”，所以优化前必须先能说清慢版本枚举了哪些候选。

\textbf{暴力过程。} 拿 A: a1->a2->c1->c2 和 B: b1->c1->c2 手算，交点是节点地址 c1，不是值相等。两个指针走到尾后切换到对方头部，会在总路程相同后于 c1 相遇。

\textbf{重复工作。} 题面模型：相交是节点地址相同，不是节点值相同。暴力版的慢点要落到本题：慢版本会围绕“相交是节点地址相同，不是节点值相同”借助数组或多次扫描确认目标顺序。样例规律是“规律是比较节点身份，长度差由切换链表自动抵消”，说明优化点不在节点值，而在少量指针角色能否保持已处理链、未处理链和接回位置都不丢。后面的优化只消掉这类重复工作，不能改变答案集合。

\textbf{优化条件。} 本题的关键观察是：两个指针走完各自链表后切换到另一条链，走过总长度相同后相遇。这句话必须能翻译成可维护状态；如果题目条件改变后观察不再成立，就不能继续套原来的写法。

\textbf{相邻状态。} 规律是比较节点身份，长度差由切换链表自动抵消。把这个特例继续拆开看：每个指针都要有角色，上一段尾、当前段头、当前节点和后继入口不能混。只要一次改 next 之前没有保存后继，就可能丢掉整段未处理链。

\textbf{状态复用。} 把数组辅助结构换成几个稳定指针角色；重连顺序必须保证已处理链合法、未处理链仍可达。本题落点是：规律是比较节点身份，长度差由切换链表自动抵消。手算时只改被新输入影响的那一小部分，而不是回头重算完整候选。

\textbf{指针和字段。} 状态字段要能说清：虚拟头、上一段尾、当前段头、当前节点、后继节点；任何改 next 的操作之前都要保存后续入口。手算时按这个协议跟踪：拿 A: a1->a2->c1->c2 和 B: b1->c1->c2 手算，交点是节点地址 c1，不是值相等。两个指针走到尾后切换到对方头部，会在总路程相同后于 c1 相遇。然后按协议复查：手算时画出 prev、cur、next 和下一段头节点。每次改 next 前先保存后继，这是链表推导的安全线。

\textbf{代码落点。} 所有重连步骤按保存后继、断开或反转、接回前缀、接回后缀的顺序写；最后检查是否形成环或丢节点。关键代码行必须能逐句翻译成“旧状态怎样变成新状态”，否则就只是模板。

\textbf{正确性。} 证明从链结构开始：已处理链连通、未处理链可达、二者连接点明确，节点没有丢失也没有成环。

\textbf{边界实验。} 先用不相交、头节点相交、尾节点相交、长度差很大做反例检查。实验时覆盖不相交、头节点相交、尾节点相交、长度差很大，每步画出前驱、当前、后继和下一段头。链表题不要只看输出值，还要检查节点是否丢失或成环。

\textbf{复杂度变化。} 暴力可能多次扫描链表或拷贝到数组；优化后每个节点通常只被访问、重连常数次，目标是线性时间和常数额外空间。

\textbf{迁移追问。} 如果题目改成“哈希集合更直观但空间更高”，先判断原状态是否仍然覆盖新答案集合，再决定是微调边界、改 value 语义，还是换算法。

\subsection{LeetCode 234 回文链表}

\textbf{题目说明。} LeetCode 234 回文链表 的题面入口要先还原成具体任务：链表回文要求前半段和后半段逆序后逐一相等。

\textbf{样例入口。} 拿 1->2->2->1 手算，快慢指针找到后半段起点，再反转后半段得到 1->2，与前半段逐个比较。

\textbf{暴力基线。} 慢版本可以先把节点顺序抄到数组里，按数组推导目标顺序。回到链表后，难点不再是值的顺序，而是节点身份和 next 指针不能丢。放到本题就是：先按“链表回文要求前半段和后半段逆序后逐一相等”完整试慢版本；样例里已经能看到“规律是链表不能随机访问，先找中点、反转后半、比较，再按需要恢复结构”，所以优化前必须先能说清慢版本枚举了哪些候选。

\textbf{暴力过程。} 拿 1->2->2->1 手算，快慢指针找到后半段起点，再反转后半段得到 1->2，与前半段逐个比较。

\textbf{重复工作。} 题面模型：链表回文要求前半段和后半段逆序后逐一相等。暴力版的慢点要落到本题：慢版本会围绕“链表回文要求前半段和后半段逆序后逐一相等”借助数组或多次扫描确认目标顺序。样例规律是“规律是链表不能随机访问，先找中点、反转后半、比较，再按需要恢复结构”，说明优化点不在节点值，而在少量指针角色能否保持已处理链、未处理链和接回位置都不丢。后面的优化只消掉这类重复工作，不能改变答案集合。

\textbf{优化条件。} 本题的关键观察是：快慢指针找中点，反转后半段，再和前半段同步比较。这句话必须能翻译成可维护状态；如果题目条件改变后观察不再成立，就不能继续套原来的写法。

\textbf{相邻状态。} 规律是链表不能随机访问，先找中点、反转后半、比较，再按需要恢复结构。把这个特例继续拆开看：每个指针都要有角色，上一段尾、当前段头、当前节点和后继入口不能混。只要一次改 next 之前没有保存后继，就可能丢掉整段未处理链。

\textbf{状态复用。} 把数组辅助结构换成几个稳定指针角色；重连顺序必须保证已处理链合法、未处理链仍可达。本题落点是：规律是链表不能随机访问，先找中点、反转后半、比较，再按需要恢复结构。手算时只改被新输入影响的那一小部分，而不是回头重算完整候选。

\textbf{指针和字段。} 状态字段要能说清：虚拟头、上一段尾、当前段头、当前节点、后继节点；任何改 next 的操作之前都要保存后续入口。手算时按这个协议跟踪：拿 1->2->2->1 手算，快慢指针找到后半段起点，再反转后半段得到 1->2，与前半段逐个比较。然后按协议复查：手算时画出 prev、cur、next 和下一段头节点。每次改 next 前先保存后继，这是链表推导的安全线。

\textbf{代码落点。} 所有重连步骤按保存后继、断开或反转、接回前缀、接回后缀的顺序写；最后检查是否形成环或丢节点。关键代码行必须能逐句翻译成“旧状态怎样变成新状态”，否则就只是模板。

\textbf{正确性。} 证明从链结构开始：已处理链连通、未处理链可达、二者连接点明确，节点没有丢失也没有成环。

\textbf{边界实验。} 先用空链、单节点、奇数长度、偶数长度、是否恢复链表做反例检查。实验时覆盖空链、单节点、奇数长度、偶数长度、是否恢复链表，每步画出前驱、当前、后继和下一段头。链表题不要只看输出值，还要检查节点是否丢失或成环。

\textbf{复杂度变化。} 暴力可能多次扫描链表或拷贝到数组；优化后每个节点通常只被访问、重连常数次，目标是线性时间和常数额外空间。

\textbf{迁移追问。} 如果题目改成“如果不允许修改链表，可用栈保存前半段但空间变为线性”，先判断原状态是否仍然覆盖新答案集合，再决定是微调边界、改 value 语义，还是换算法。

\subsection{LeetCode 148 排序链表}

\textbf{题目说明。} LeetCode 148 排序链表 的题面入口要先还原成具体任务：链表不能随机访问，适合归并排序而不是快速排序数组 partition。

\textbf{样例入口。} 拿 4->2->1->3 手算，快慢指针把链表切成 4->2 和 1->3，分别排好后再归并。

\textbf{暴力基线。} 慢版本可以先把节点顺序抄到数组里，按数组推导目标顺序。回到链表后，难点不再是值的顺序，而是节点身份和 next 指针不能丢。放到本题就是：先按“链表不能随机访问，适合归并排序而不是快速排序数组 partition”完整试慢版本；样例里已经能看到“规律是链表排序不能随机访问，归并排序只需要切断中点和有序合并，复杂度稳定在 O(n log n)”，所以优化前必须先能说清慢版本枚举了哪些候选。

\textbf{暴力过程。} 拿 4->2->1->3 手算，快慢指针把链表切成 4->2 和 1->3，分别排好后再归并。

\textbf{重复工作。} 题面模型：链表不能随机访问，适合归并排序而不是快速排序数组 partition。暴力版的慢点要落到本题：慢版本会围绕“链表不能随机访问，适合归并排序而不是快速排序数组 partition”借助数组或多次扫描确认目标顺序。样例规律是“规律是链表排序不能随机访问，归并排序只需要切断中点和有序合并，复杂度稳定在 O(n log n)”，说明优化点不在节点值，而在少量指针角色能否保持已处理链、未处理链和接回位置都不丢。后面的优化只消掉这类重复工作，不能改变答案集合。

\textbf{优化条件。} 本题的关键观察是：自底向上迭代归并能避免递归栈，同时保持 O(n log n)。这句话必须能翻译成可维护状态；如果题目条件改变后观察不再成立，就不能继续套原来的写法。

\textbf{相邻状态。} 规律是链表排序不能随机访问，归并排序只需要切断中点和有序合并，复杂度稳定在 O(n log n)。把这个特例继续拆开看：每个指针都要有角色，上一段尾、当前段头、当前节点和后继入口不能混。只要一次改 next 之前没有保存后继，就可能丢掉整段未处理链。

\textbf{状态复用。} 把数组辅助结构换成几个稳定指针角色；重连顺序必须保证已处理链合法、未处理链仍可达。本题落点是：规律是链表排序不能随机访问，归并排序只需要切断中点和有序合并，复杂度稳定在 O(n log n)。手算时只改被新输入影响的那一小部分，而不是回头重算完整候选。

\textbf{指针和字段。} 状态字段要能说清：虚拟头、上一段尾、当前段头、当前节点、后继节点；任何改 next 的操作之前都要保存后续入口。手算时按这个协议跟踪：拿 4->2->1->3 手算，快慢指针把链表切成 4->2 和 1->3，分别排好后再归并。然后按协议复查：手算时画出 prev、cur、next 和下一段头节点。每次改 next 前先保存后继，这是链表推导的安全线。

\textbf{代码落点。} 所有重连步骤按保存后继、断开或反转、接回前缀、接回后缀的顺序写；最后检查是否形成环或丢节点。关键代码行必须能逐句翻译成“旧状态怎样变成新状态”，否则就只是模板。

\textbf{正确性。} 证明从链结构开始：已处理链连通、未处理链可达、二者连接点明确，节点没有丢失也没有成环。

\textbf{边界实验。} 先用空链、单节点、重复值、奇数长度、切断子链做反例检查。实验时覆盖空链、单节点、重复值、奇数长度、切断子链，每步画出前驱、当前、后继和下一段头。链表题不要只看输出值，还要检查节点是否丢失或成环。

\textbf{复杂度变化。} 暴力可能多次扫描链表或拷贝到数组；优化后每个节点通常只被访问、重连常数次，目标是线性时间和常数额外空间。

\textbf{迁移追问。} 如果题目改成“若把值复制到数组排序，会改变节点身份语义”，先判断原状态是否仍然覆盖新答案集合，再决定是微调边界、改 value 语义，还是换算法。

\subsection{LeetCode 82 删除排序链表中的重复元素 II}

\textbf{题目说明。} LeetCode 82 删除排序链表中的重复元素 II 的题面入口要先还原成具体任务：有序链表中重复值连续出现，题目要求把所有重复值整组删除。

\textbf{样例入口。} 拿 1->2->3->3->4 手算，两个 3 都要删除，不是保留一个；pre 停在 2，cur 扫完整段重复 3 后接到 4。

\textbf{暴力基线。} 慢版本可以先把节点顺序抄到数组里，按数组推导目标顺序。回到链表后，难点不再是值的顺序，而是节点身份和 next 指针不能丢。放到本题就是：先按“有序链表中重复值连续出现，题目要求把所有重复值整组删除”完整试慢版本；样例里已经能看到“规律是虚拟头加 pre 指向已确认无重复前缀尾，遇到重复段时整段跳过”，所以优化前必须先能说清慢版本枚举了哪些候选。

\textbf{暴力过程。} 拿 1->2->3->3->4 手算，两个 3 都要删除，不是保留一个；pre 停在 2，cur 扫完整段重复 3 后接到 4。

\textbf{重复工作。} 题面模型：有序链表中重复值连续出现，题目要求把所有重复值整组删除。暴力版的慢点要落到本题：慢版本会围绕“有序链表中重复值连续出现，题目要求把所有重复值整组删除”借助数组或多次扫描确认目标顺序。样例规律是“规律是虚拟头加 pre 指向已确认无重复前缀尾，遇到重复段时整段跳过”，说明优化点不在节点值，而在少量指针角色能否保持已处理链、未处理链和接回位置都不丢。后面的优化只消掉这类重复工作，不能改变答案集合。

\textbf{优化条件。} 本题的关键观察是：虚拟头保存结果前驱，遇到一段重复值时跳过整段，否则前驱前进。这句话必须能翻译成可维护状态；如果题目条件改变后观察不再成立，就不能继续套原来的写法。

\textbf{相邻状态。} 规律是虚拟头加 pre 指向已确认无重复前缀尾，遇到重复段时整段跳过。把这个特例继续拆开看：每个指针都要有角色，上一段尾、当前段头、当前节点和后继入口不能混。只要一次改 next 之前没有保存后继，就可能丢掉整段未处理链。

\textbf{状态复用。} 把数组辅助结构换成几个稳定指针角色；重连顺序必须保证已处理链合法、未处理链仍可达。本题落点是：规律是虚拟头加 pre 指向已确认无重复前缀尾，遇到重复段时整段跳过。手算时只改被新输入影响的那一小部分，而不是回头重算完整候选。

\textbf{指针和字段。} 状态字段要能说清：虚拟头、上一段尾、当前段头、当前节点、后继节点；任何改 next 的操作之前都要保存后续入口。手算时按这个协议跟踪：拿 1->2->3->3->4 手算，两个 3 都要删除，不是保留一个；pre 停在 2，cur 扫完整段重复 3 后接到 4。然后按协议复查：手算时画出 prev、cur、next 和下一段头节点。每次改 next 前先保存后继，这是链表推导的安全线。

\textbf{代码落点。} 所有重连步骤按保存后继、断开或反转、接回前缀、接回后缀的顺序写；最后检查是否形成环或丢节点。关键代码行必须能逐句翻译成“旧状态怎样变成新状态”，否则就只是模板。

\textbf{正确性。} 证明从链结构开始：已处理链连通、未处理链可达、二者连接点明确，节点没有丢失也没有成环。

\textbf{边界实验。} 先用头部重复、尾部重复、全部重复、没有重复做反例检查。实验时覆盖头部重复、尾部重复、全部重复、没有重复，每步画出前驱、当前、后继和下一段头。链表题不要只看输出值，还要检查节点是否丢失或成环。

\textbf{复杂度变化。} 暴力可能多次扫描链表或拷贝到数组；优化后每个节点通常只被访问、重连常数次，目标是线性时间和常数额外空间。

\textbf{迁移追问。} 如果题目改成“若只保留一个重复值，就是删除重复元素的较简单版本”，先判断原状态是否仍然覆盖新答案集合，再决定是微调边界、改 value 语义，还是换算法。

\subsection{LeetCode 71 简化路径}

\textbf{题目说明。} LeetCode 71 简化路径 的题面入口要先还原成具体任务：Unix 绝对路径由斜杠分隔目录段，点号段和返回上级目录段需要被规范化。

\textbf{样例入口。} 拿 path=/home//foo/ 手算，连续斜杠产生空段要忽略，home 和 foo 依次入栈，结果是 /home/foo；拿 /../ 手算，根目录上遇到 .. 不能再向上弹。

\textbf{暴力基线。} 慢版本让每个位置向左或向右线性寻找边界，或者让每个结构反复等待匹配。它能算出答案，但很多尚未完成的候选会被未来同一个事件一起解决。放到本题就是：先按“Unix 绝对路径由斜杠分隔目录段，点号段和返回上级目录段需要被规范化”完整试慢版本；样例里已经能看到“规律是按斜杠切分段，普通目录入栈，.. 弹出有效目录，空段和 . 忽略”，所以优化前必须先能说清慢版本枚举了哪些候选。

\textbf{暴力过程。} 拿 path=/home//foo/ 手算，连续斜杠产生空段要忽略，home 和 foo 依次入栈，结果是 /home/foo；拿 /../ 手算，根目录上遇到 .. 不能再向上弹。

\textbf{重复工作。} 题面模型：Unix 绝对路径由斜杠分隔目录段，点号段和返回上级目录段需要被规范化。暴力版的慢点要落到本题：慢版本会围绕“Unix 绝对路径由斜杠分隔目录段，点号段和返回上级目录段需要被规范化”反复寻找最近边界或最近未完成对象。样例规律是“规律是按斜杠切分段，普通目录入栈，.. 弹出有效目录，空段和 . 忽略”，说明旧候选一旦遇到当前输入就能被批量结算；继续为每个位置单独向左或向右扫描就是浪费。后面的优化只消掉这类重复工作，不能改变答案集合。

\textbf{优化条件。} 本题的关键观察是：按斜杠切分路径，普通目录入栈，遇到 .. 弹出一个有效目录，空段和 . 忽略。这句话必须能翻译成可维护状态；如果题目条件改变后观察不再成立，就不能继续套原来的写法。

\textbf{相邻状态。} 规律是按斜杠切分段，普通目录入栈，.. 弹出有效目录，空段和 . 忽略。把这个特例继续拆开看：栈里的对象都是答案尚未确定的候选；一旦当前元素触发弹栈，被弹出的候选必须已经同时拿到了左边界和右边界，未来元素不再有资格改变它的答案。

\textbf{状态复用。} 把反复寻找最近边界改成维护未完成候选；新元素只和栈顶比较，连续弹出一批已经被当前元素解决的对象。本题落点是：规律是按斜杠切分段，普通目录入栈，.. 弹出有效目录，空段和 . 忽略。手算时只改被新输入影响的那一小部分，而不是回头重算完整候选。

\textbf{指针和字段。} 状态字段要能说清：栈内对象的业务含义、当前输入、弹出时结算的对象、左边界和右边界；栈中存下标还是值要明确。手算时按这个协议跟踪：拿 path=/home//foo/ 手算，连续斜杠产生空段要忽略，home 和 foo 依次入栈，结果是 /home/foo；拿 /../ 手算，根目录上遇到 .. 不能再向上弹。然后按协议复查：手算时记录栈内元素的业务含义，不只记录数值。入栈表示答案未定，弹栈表示当前事件已经让它的答案定下来。

\textbf{代码落点。} 弹栈前确认栈非空，弹出后再读取新的左边界；相等元素和哨兵高度要有统一策略。关键代码行必须能逐句翻译成“旧状态怎样变成新状态”，否则就只是模板。

\textbf{正确性。} 证明从弹出时机开始：被弹出的对象在当前时刻答案已经确定，未来元素无法再改变它。

\textbf{边界实验。} 先用连续斜杠、根目录上继续 ..、路径末尾斜杠、目录名含点但不是点号做反例检查。实验时准备连续斜杠、根目录上继续 ..、路径末尾斜杠、目录名含点但不是点号，打印每次入栈、弹栈和结算对象。重点观察相等元素、空栈、哨兵和最后一次结算。

\textbf{复杂度变化。} 暴力为每个位置寻找边界常见为平方级；单调栈让每个元素最多入栈一次、出栈一次，因此主体扫描为线性时间。

\textbf{迁移追问。} 如果题目改成“若路径是相对路径，开头和越过根目录的处理规则要重新定义”，先判断原状态是否仍然覆盖新答案集合，再决定是微调边界、改 value 语义，还是换算法。

\subsection{LeetCode 32 最长有效括号}

\textbf{题目说明。} LeetCode 32 最长有效括号 的题面入口要先还原成具体任务：有效括号子串需要知道每段非法边界之后的最早可连接位置。

\textbf{样例入口。} 拿字符串 )()()) 手算。第一个 ) 不是可匹配对象，而是非法边界；后面每次匹配成功，当前有效长度由当前位置减去栈顶边界得出。若栈空，要把当前右括号作为新非法边界压入。

\textbf{暴力基线。} 慢版本让每个位置向左或向右线性寻找边界，或者让每个结构反复等待匹配。它能算出答案，但很多尚未完成的候选会被未来同一个事件一起解决。放到本题就是：先按“有效括号子串需要知道每段非法边界之后的最早可连接位置”完整试慢版本；样例里已经能看到“规律是栈里既保存未匹配左括号下标，也保存最近非法边界”，所以优化前必须先能说清慢版本枚举了哪些候选。

\textbf{暴力过程。} 拿字符串 )()()) 手算。第一个 ) 不是可匹配对象，而是非法边界；后面每次匹配成功，当前有效长度由当前位置减去栈顶边界得出。若栈空，要把当前右括号作为新非法边界压入。

\textbf{重复工作。} 题面模型：有效括号子串需要知道每段非法边界之后的最早可连接位置。暴力版的慢点要落到本题：慢版本会围绕“有效括号子串需要知道每段非法边界之后的最早可连接位置”反复寻找最近边界或最近未完成对象。样例规律是“规律是栈里既保存未匹配左括号下标，也保存最近非法边界”，说明旧候选一旦遇到当前输入就能被批量结算；继续为每个位置单独向左或向右扫描就是浪费。后面的优化只消掉这类重复工作，不能改变答案集合。

\textbf{优化条件。} 本题的关键观察是：栈保存下标，先放入哨兵负一；匹配后用当前下标减栈顶得到长度。这句话必须能翻译成可维护状态；如果题目条件改变后观察不再成立，就不能继续套原来的写法。

\textbf{相邻状态。} 规律是栈里既保存未匹配左括号下标，也保存最近非法边界。把这个特例继续拆开看：栈里的对象都是答案尚未确定的候选；一旦当前元素触发弹栈，被弹出的候选必须已经同时拿到了左边界和右边界，未来元素不再有资格改变它的答案。

\textbf{状态复用。} 把反复寻找最近边界改成维护未完成候选；新元素只和栈顶比较，连续弹出一批已经被当前元素解决的对象。本题落点是：规律是栈里既保存未匹配左括号下标，也保存最近非法边界。手算时只改被新输入影响的那一小部分，而不是回头重算完整候选。

\textbf{指针和字段。} 状态字段要能说清：栈内对象的业务含义、当前输入、弹出时结算的对象、左边界和右边界；栈中存下标还是值要明确。手算时按这个协议跟踪：拿字符串 )()()) 手算。第一个 ) 不是可匹配对象，而是非法边界；后面每次匹配成功，当前有效长度由当前位置减去栈顶边界得出。若栈空，要把当前右括号作为新非法边界压入。然后按协议复查：手算时记录栈内元素的业务含义，不只记录数值。入栈表示答案未定，弹栈表示当前事件已经让它的答案定下来。

\textbf{代码落点。} 弹栈前确认栈非空，弹出后再读取新的左边界；相等元素和哨兵高度要有统一策略。关键代码行必须能逐句翻译成“旧状态怎样变成新状态”，否则就只是模板。

\textbf{正确性。} 证明从弹出时机开始：被弹出的对象在当前时刻答案已经确定，未来元素无法再改变它。

\textbf{边界实验。} 先用空串、全左括号、全右括号、多个断开的合法段做反例检查。实验时准备空串、全左括号、全右括号、多个断开的合法段，打印每次入栈、弹栈和结算对象。重点观察相等元素、空栈、哨兵和最后一次结算。

\textbf{复杂度变化。} 暴力为每个位置寻找边界常见为平方级；单调栈让每个元素最多入栈一次、出栈一次，因此主体扫描为线性时间。

\textbf{迁移追问。} 如果题目改成“DP 写法也可做，但栈写法更直接表达最近非法边界”，先判断原状态是否仍然覆盖新答案集合，再决定是微调边界、改 value 语义，还是换算法。

\subsection{LeetCode 206 反转链表}

\textbf{题目说明。} LeetCode 206 反转链表 的题面入口要先还原成具体任务：每次把当前节点的 next 指向已经反转好的前缀头。

\textbf{样例入口。} 拿 1->2->3 手算，处理 1 时先保存 next=2，再让 1->null；处理 2 时保存 3，再让 2->1。

\textbf{暴力基线。} 慢版本可以先把节点顺序抄到数组里，按数组推导目标顺序。回到链表后，难点不再是值的顺序，而是节点身份和 next 指针不能丢。放到本题就是：先按“每次把当前节点的 next 指向已经反转好的前缀头”完整试慢版本；样例里已经能看到“规律是 prev 保存已反转前缀头，cur 保存待处理节点，每次改 next 前必须保存后继”，所以优化前必须先能说清慢版本枚举了哪些候选。

\textbf{暴力过程。} 拿 1->2->3 手算，处理 1 时先保存 next=2，再让 1->null；处理 2 时保存 3，再让 2->1。

\textbf{重复工作。} 题面模型：每次把当前节点的 next 指向已经反转好的前缀头。暴力版的慢点要落到本题：慢版本会围绕“每次把当前节点的 next 指向已经反转好的前缀头”借助数组或多次扫描确认目标顺序。样例规律是“规律是 prev 保存已反转前缀头，cur 保存待处理节点，每次改 next 前必须保存后继”，说明优化点不在节点值，而在少量指针角色能否保持已处理链、未处理链和接回位置都不丢。后面的优化只消掉这类重复工作，不能改变答案集合。

\textbf{优化条件。} 本题的关键观察是：先保存后继，再改指针，再推进 previous 和 current。这句话必须能翻译成可维护状态；如果题目条件改变后观察不再成立，就不能继续套原来的写法。

\textbf{相邻状态。} 规律是 prev 保存已反转前缀头，cur 保存待处理节点，每次改 next 前必须保存后继。把这个特例继续拆开看：每个指针都要有角色，上一段尾、当前段头、当前节点和后继入口不能混。只要一次改 next 之前没有保存后继，就可能丢掉整段未处理链。

\textbf{状态复用。} 把数组辅助结构换成几个稳定指针角色；重连顺序必须保证已处理链合法、未处理链仍可达。本题落点是：规律是 prev 保存已反转前缀头，cur 保存待处理节点，每次改 next 前必须保存后继。手算时只改被新输入影响的那一小部分，而不是回头重算完整候选。

\textbf{指针和字段。} 状态字段要能说清：虚拟头、上一段尾、当前段头、当前节点、后继节点；任何改 next 的操作之前都要保存后续入口。手算时按这个协议跟踪：拿 1->2->3 手算，处理 1 时先保存 next=2，再让 1->null；处理 2 时保存 3，再让 2->1。然后按协议复查：手算时画出 prev、cur、next 和下一段头节点。每次改 next 前先保存后继，这是链表推导的安全线。

\textbf{代码落点。} 所有重连步骤按保存后继、断开或反转、接回前缀、接回后缀的顺序写；最后检查是否形成环或丢节点。关键代码行必须能逐句翻译成“旧状态怎样变成新状态”，否则就只是模板。

\textbf{正确性。} 证明从链结构开始：已处理链连通、未处理链可达、二者连接点明确，节点没有丢失也没有成环。

\textbf{边界实验。} 先用空链、单节点、两个节点、不要丢失剩余链做反例检查。实验时覆盖空链、单节点、两个节点、不要丢失剩余链，每步画出前驱、当前、后继和下一段头。链表题不要只看输出值，还要检查节点是否丢失或成环。

\textbf{复杂度变化。} 暴力可能多次扫描链表或拷贝到数组；优化后每个节点通常只被访问、重连常数次，目标是线性时间和常数额外空间。

\textbf{迁移追问。} 如果题目改成“K 组反转是在这段局部反转外增加分组边界”，先判断原状态是否仍然覆盖新答案集合，再决定是微调边界、改 value 语义，还是换算法。

\subsection{LeetCode 739 每日温度}

\textbf{题目说明。} LeetCode 739 每日温度 的题面入口要先还原成具体任务：每一天等待右侧第一个更高温度。

\textbf{样例入口。} 拿 [73, 74, 75, 71, 69, 72, 76, 73] 手算，73 遇到 74 时等待 1 天；71 要等到 72 才解决。

\textbf{暴力基线。} 慢版本让每个位置向左或向右线性寻找边界，或者让每个结构反复等待匹配。它能算出答案，但很多尚未完成的候选会被未来同一个事件一起解决。放到本题就是：先按“每一天等待右侧第一个更高温度”完整试慢版本；样例里已经能看到“规律是单调递减栈保存尚未等到更高温的下标，遇到更高温时连续弹栈结算天数”，所以优化前必须先能说清慢版本枚举了哪些候选。

\textbf{暴力过程。} 拿 [73, 74, 75, 71, 69, 72, 76, 73] 手算，73 遇到 74 时等待 1 天；71 要等到 72 才解决。

\textbf{重复工作。} 题面模型：每一天等待右侧第一个更高温度。暴力版的慢点要落到本题：慢版本会围绕“每一天等待右侧第一个更高温度”反复寻找最近边界或最近未完成对象。样例规律是“规律是单调递减栈保存尚未等到更高温的下标，遇到更高温时连续弹栈结算天数”，说明旧候选一旦遇到当前输入就能被批量结算；继续为每个位置单独向左或向右扫描就是浪费。后面的优化只消掉这类重复工作，不能改变答案集合。

\textbf{优化条件。} 本题的关键观察是：单调递减栈保存尚未找到答案的下标。这句话必须能翻译成可维护状态；如果题目条件改变后观察不再成立，就不能继续套原来的写法。

\textbf{相邻状态。} 规律是单调递减栈保存尚未等到更高温的下标，遇到更高温时连续弹栈结算天数。把这个特例继续拆开看：栈里的对象都是答案尚未确定的候选；一旦当前元素触发弹栈，被弹出的候选必须已经同时拿到了左边界和右边界，未来元素不再有资格改变它的答案。

\textbf{状态复用。} 把反复寻找最近边界改成维护未完成候选；新元素只和栈顶比较，连续弹出一批已经被当前元素解决的对象。本题落点是：规律是单调递减栈保存尚未等到更高温的下标，遇到更高温时连续弹栈结算天数。手算时只改被新输入影响的那一小部分，而不是回头重算完整候选。

\textbf{指针和字段。} 状态字段要能说清：栈内对象的业务含义、当前输入、弹出时结算的对象、左边界和右边界；栈中存下标还是值要明确。手算时按这个协议跟踪：拿 [73, 74, 75, 71, 69, 72, 76, 73] 手算，73 遇到 74 时等待 1 天；71 要等到 72 才解决。然后按协议复查：手算时记录栈内元素的业务含义，不只记录数值。入栈表示答案未定，弹栈表示当前事件已经让它的答案定下来。

\textbf{代码落点。} 弹栈前确认栈非空，弹出后再读取新的左边界；相等元素和哨兵高度要有统一策略。关键代码行必须能逐句翻译成“旧状态怎样变成新状态”，否则就只是模板。

\textbf{正确性。} 证明从弹出时机开始：被弹出的对象在当前时刻答案已经确定，未来元素无法再改变它。

\textbf{边界实验。} 先用没有更高温、相等温度、递增、递减做反例检查。实验时准备没有更高温、相等温度、递增、递减，打印每次入栈、弹栈和结算对象。重点观察相等元素、空栈、哨兵和最后一次结算。

\textbf{复杂度变化。} 暴力为每个位置寻找边界常见为平方级；单调栈让每个元素最多入栈一次、出栈一次，因此主体扫描为线性时间。

\textbf{迁移追问。} 如果题目改成“下一个更大元素与它同型，只是输出值或距离不同”，先判断原状态是否仍然覆盖新答案集合，再决定是微调边界、改 value 语义，还是换算法。

\begin{principlebox}{本节复盘方式}
不要只看最终算法名。每道题复盘时先遮住优化版，口述暴力枚举对象；再指出相邻窗口、历史前缀、候选区间、搜索状态或子问题里哪一部分被重复处理；最后用一个边界样例验证状态更新顺序。能讲完这三步，才算真正掌握本章案例。
\end{principlebox}
